\subsection{时间的涌现：从状态精化到时间投影}\label{subsec:physical-spacetime-skeleton-time-projection}

上一小节已经把时空的最小骨架压缩到事件、区域与覆盖，但在那里，时间仍只以一个尚待解释的位置出现：事件集 \(X\) 上虽带有单调的时间索引，状态也可经由事件定位映射落到某个事件上，然而这还不足以说明时间本身究竟是什么。若在此处直接把时间索引读作一个独立于状态变化的背景参数，那么整个体系就会重新退回到“先有时间容器，再有状态演化”的旧写法，与本章一开始所坚持的立场相冲突。因此，接下来的关键并不是为系统再加上一层时间公理，而是说明：在统一上位逻辑中，时间本来就不是前提，而是状态精化、残差外置与事件定位共同诱导出的投影量。

要看清这一点，首先必须回到状态更新本身。\S\ref{subsec:logic-expansion-chain-dynamics-time-multiaxis} 已经把更新写成状态的进一步细化，而不是对既有事实的回写；定义 \ref{def:logic-expansion-multi-axis-dynamic-structure} 与命题 \ref{prop:logic-expansion-updates-preserve-forcing} 表明，一个更新并不是把旧状态整体废弃后再任意换成一个新状态，而是把原本较粗的局部信息推进到更细的状态，使先前尚未排除的可能性被进一步剔除，使先前尚未分辨的分支被进一步分开。在这一意义上，更新链天然带有先后：不是因为系统先被放进了一条时间轴，而是因为“较细”总在“较粗”之后，“已分辨”总在“未分辨”之后。于是，最原始的时间读法首先不是钟表式的长度，而是状态精化所形成的顺序结构。谁先谁后，不是由外部背景授予，而是由细化关系本身决定。

一旦这一点成立，事件时间的地位也就必须随之改写。对每个状态 \(p\)，定义 \ref{def:logic-expansion-observer-spacetime-state-model} 已经给出其事件定位 \(\lambda(p)\in X\)，而事件再经由单调映射投向时间索引；相应时间投影写作
\[
\tau=\chi\circ\lambda.
\]
于是，状态的“时间坐标”并不应被理解为某个独立背景变量，而应理解为组合映射 \(\chi\circ\lambda\) 在状态上的取值。换言之，状态并不是先在一个外部时间中流动，然后再被附着到某个事件上；恰恰相反，正是因为状态已经通过更新关系形成了可排序的链，并且每个状态都被定位到了某个事件位置上，我们才可以把这个顺序压缩成一个可读的时间投影。于是，时间不是与状态流动并列的第二条轴，而是状态更新秩序在事件层上的可视化结果。

不过，若只把时间理解为更新先后，仍然还不够。因为前文真正具有辨识度的并不只是一般的状态更新，而是折叠、局部读出与纤维残差同时存在时，那种“对象层已定而分辨尚未完成”的情形。折叠之后，多个微观状态可以投到同一可见对象上；从对象层看，读出似乎已经完成，但从实现层看，仍有一部分纤维自由度没有被吸收。这部分未被读出的剩余，在 \S\ref{subsec:residual-unfolding-time-depth} 中已经被系统地写成残差，并指出它至少有两种等价口径：一方面，它表现为实现层尚需补入的侧信息预算；另一方面，它又表现为对象层上不可见但在纤维内部仍真实存在的分支差异。由此可见，在这套体系里，真正推动“时间化”发生的，并不只是一般意义上的更新，而是残差何时、以何种方式、在多深的层级上被进一步展开。

这正是“时间作为解析深度”这一读法的核心。\S\ref{subsubsec:frt-time-as-depth} 已经把前缀超度量、分辨率分离与折叠分离深度 \(m_{\mathrm{sep}}\) 组织为同一层级结构：两个状态究竟在多深的层级上才开始分叉，不再只是直观说法，而成为可度量的解析事实。共享前缀越长，说明它们在较粗分辨率下越难区分；必须进入更深的层级，更多地展开此前被折叠到纤维中的信息，才能把它们分开。于是，时间的一个更内在的定义就出现了：时间不是某个先验流逝量，而是为了把残差转化为可断言差异所必需的最小解析深度。所谓“更晚才被看清”，其精确含义并不是它晚发生，而是要在更深的分辨率层级上，系统才有足够的信息把它与其他候选状态区分开来。

从这里再往前一步，前文关于“序列层可逆化”与“不可逆性外置为记录轴”的结果就变得尤为重要。因为只要对象层的读出本身不足以保留全部纤维残差，而我们又要求扩展保持与可审计可逆性，那么系统就不能把这部分残差简单抹去，只能把它外置到一个附加的记录方向上。\S\ref{subsubsec:frt-sequence-inversion} 已经表明，窗口层被折叠压缩掉的纤维自由度，在阈值后并未消失，而是被系统性迁移为序列层逆码 \(\Psi_m\) 所需的有限记忆深度；与之对应，定义 \ref{def:pom-value-delay-rewrite} 和注 \ref{rem:pom-value-delay-normal-form} 把对象层不可见演化段中的值投影推迟到末端一次执行，进一步说明“延后显现”并不是事实迟到生成，而是不可见残差被按规则推迟到记录轴上结算。正是在这一点上，时间不再只是静态顺序，而成为一种外置记录的方向：对象层之所以显得“向前流动”，并不是因为世界被放在一条绝对时间线上，而是因为残差不能同时在当前层被全部消化，只能沿着记录轴按顺序展开。

这样一来，最小在线延迟便自然出现了。若残差的外置不是任意选择，而是由可逆化纪律逼出的必要步骤，那么在读出与完全恢复之间，就会存在一个不可再压缩的最小步数。\S\ref{subsubsec:frt-min-online-delay} 已将这一点写成可审计的机制下界：定理 \ref{thm:online-delay-fold} 给出一个延迟 \(3\) 的 \(\Fold_\infty\) 单趟在线实现，而命题 \ref{prop:pom-delay-min} 进一步证明该常数是最小的。这个步数不是心理上的等待时间，也不是几何上的钟表时长，而是系统为使残差从“外置但未解码”转入“可恢复且可核验”所必须支付的最小在线更新数。也正是在这里，时间第一次获得了可操作的量化含义：它不再只是抽象顺序，而是最小恢复延迟。若借用操作层语言，则所谓“时间量子”在本体系中的首要含义并不是自然界先有一格一格的时钟，而是残差外化过程存在不可约的最小步长，因此时间在操作层面呈现出离散的、不可无限压缩的粒度。

由此可以看到，本体系中的时间实际上有三层彼此嵌套的读法，而且这三层并不是三套不同理论，而是同一结构从不同角度的展开。最弱的一层是逻辑时间，即状态精化链上的先后顺序；更强的一层是解析时间，即为了分辨纤维残差所需的最小分辨深度；最强的一层则是操作时间，即残差外置后实现可恢复所必须经历的最小在线延迟。三者的共同点在于：它们都不是外加背景，而都是状态—残差结构自身的不同投影。也正因此，本章完全不必为时间另立公理；只要承认前文已经建立的细化、折叠、残差与可逆化纪律，时间就已经以这三重方式同时出现了。

在这种理解下，时间的物理化读法也获得了一个非常重要的边界。若时间只是状态精化与残差展开的投影，那么“过去”就不应再被理解为一个可以被后续操作重新改写的对象层库存；后续更新真正改变的，只是较早命题在当前状态中的可断言性，以及相关残差是否已经被排除或解码。换言之，这一投影性读法天然把“后验分辨”与“回写过去”区分开来：前者是细化链上的定解，后者则会违反命题 \ref{prop:logic-expansion-delayed-decidability-no-new-truth} 与命题 \ref{prop:observer-indexed-delayed-classification-not-retrocausal} 所固定的不回写纪律。下一小节便将在这一基础上进一步说明：延迟观测之所以可能，并不意味着过去对象层内容被重新书写，而只意味着先前未定的局部可断言性，在后续状态中得到了排除与定解。
